% !TEX root = ../main.tex

\chapter{Android E4 Reproducibility}
\label{app:android-e4}

\section{Environment and Protocol}

The Android E4 measurements used an Android 14/API 34 \texttt{sdk\_gphone64\_x86\_64} emulator with four virtual processors and a 192~MiB application heap.
The live-backend trace used the emulator's host-network route to a local backend.
The authentication token was supplied to the instrumentation process at run time and was not retained in an artifact.
The retained manifest binds these measurements to the evaluated Android source snapshot.
This snapshot bounds the client results in Section~\ref{sec:e4-android}; it does not change the evidence status of the other evaluation tracks.

Startup measurement discarded three cold/foreground launch pairs, then retained 20 cold process starts and 20 task foreground/resume operations.
All retained operations returned a non-empty \texttt{WaitTime} field for the requested activity; the startup runner did not inspect rendered UI content.
Android classified nine as hot starts and returned \texttt{UNKNOWN (0)} for eleven, so their latency distribution uses the retained wait-time field and is not treated as a 20-sample confirmed hot-start distribution.
Controlled state, graph, and event tests used separate warm-up iterations that were not included in the reported samples.
The graph experiment used deterministic payloads with 1, 12, 25, and 50 nodes.
The live event loop performed five repetitions.

\section{Metric Definitions}

\begin{longtable}{P{3.7cm}P{10.4cm}}
\caption{Android E4 metric definitions.}
\label{tab:e4-metric-definitions}\\
\toprule
Metric & Definition \\
\midrule
Cold process start &
Elapsed time from a forced process start to the expected first Android screen. \\
\addlinespace
Task foreground/resume &
Elapsed wait time from requesting the existing task in the foreground.
Success requires a non-empty activity wait-time record; Android's hot-start classifier is recorded separately. \\
\addlinespace
Controlled state transition &
Elapsed time from delivering a deterministic fixture result to the expected ViewModel state.
This excludes live network time and display-frame latency. \\
\addlinespace
Three-stage polling &
Elapsed time until a controlled asynchronous job advances through three client-visible stages and reaches ready state under the production one-second polling schedule. \\
\addlinespace
Graph DOM-ready &
Elapsed time from loading or updating the local graph page to its JavaScript readiness callback.
The metric does not wait for force-layout convergence. \\
\addlinespace
Selection callback &
Elapsed time from selecting a graph node to receipt of the validated identifier in the native Android layer. \\
\addlinespace
Event stage latency &
Elapsed time for the named queue, retry, replay, upload, or readback operation.
Controlled and live stages are summarized separately. \\
\addlinespace
Median and p95 &
The median is the sample median.
The p95 is the nearest-rank 95th percentile, so the maximum is reported for a five-sample live stage. \\
\bottomrule
\end{longtable}

\section{Sample Inventory}

\begin{longtable}{P{4.2cm}P{1.8cm}P{2.7cm}P{3.7cm}}
\caption{Retained sample inventory for the Android E4 evaluation.}
\label{tab:e4-sample-inventory}\\
\toprule
Group & Samples & Conditions met & Included conditions \\
\midrule
Startup and resume & 40 & 40 & 20 cold starts; 20 foreground/resume operations \\
State and polling & 160 & 160 & five state conditions at 30 trials; ten polling trials \\
Graph timing & 400 & 400 & ready and selection; four sizes; new and existing WebView \\
Graph-state disclosure & 90 & 90 & ready, fallback, and empty; 30 trials each \\
Controlled event sync & 120 & 120 & queue, failure, retry, and replay; 30 trials each \\
\midrule
Controlled total & 810 & 810 & five result groups reported separately in Chapter~4 \\
\addlinespace
Live-backend stages & 25 & 25 & five stages in each of five repetitions \\
Existing graph/app-flow tests & 11 & 11 & includes Open Paper and restored selection \\
\bottomrule
\end{longtable}

\section{Success Conditions and Retained Artifacts}

A cold or foreground/resume launch succeeds when \texttt{am start -W} returns a non-empty wait-time record for the requested activity.
This condition does not inspect rendered UI content.
A content-state trial succeeds when the ViewModel exposes the expected content origin or retryable error.
A polling trial succeeds when all three processing stages are observed in order and the paper reaches ready state.
A graph trial succeeds when the page reports ready and the selected identifier reaches the native client.
Ready, fallback, and empty tests additionally require the corresponding reader-visible disclosure.

An event trial succeeds when the local queue and reservation state match the expected stage.
After a retryable failure, the batch must return to pending state.
After success, accepted and duplicate counts must account for the full batch.
For the live trace, each repetition must queue three events, recover from the injected failure, upload three new identifiers, recognize all three on replay, and return a live-origin briefing to the client.
Digest entry count and preference-model version are recorded as outcomes rather than success conditions.

The retained evaluation package contains direct comma-separated measurements, an environment record, a run manifest without the authentication token, connected Android test XML, machine-readable summaries, and the vector figures used in Chapter~4.
The five live digests contain zero entries and retain preference-model version 1.
Those values remain part of the reproducibility record because they bound the event-loop conclusion.
